\begin{questionS}
Consider the following scenario.  There are two types of threads, type A and
type B, that share a resource.  Multiple threads of the same type may use the
resource simultaneously.  However, threads of different types may not.

We want a lock object to control access in this scenario.  The lock should
have the following interface.
%
\begin{scala}
/** Trait for a disjoint usage lock. */
trait DisjointLock{
  /** An A-thread enters the critical section. */
  def aEnter(): Unit
  /** An A-thread leaves the critical section. */
  def aExit(): Unit
  /** A B-thread enters the critical section. */
  def bEnter(): Unit
  /** A B-thread leaves the critical section. */
  def bExit(): Unit
}
\end{scala}
%
The |aEnter| and |bEnter| operations should block while a thread of the other
type is using the resource.  

Implement a class with this interface, using conditions internally.  Your
implementation should avoid starvation: assuming no single thread uses the
resource forever, threads of the other type should eventually gain access. 

Test your implementation using linearizability testing. 
\end{questionS}

%%%%%%%%%%%%%%%%%%%%%%%%%%%%%%%%%%%%%%%%%%%%%%%%%%%%%%%

\begin{answerS}
My code is below, explained via comments.  This adapts the tactic from the
readers and writers example: a new thread does not enter the critical region
immediately if there are threads of the other type waiting. 
%
\begin{scala}
class MonitorDisjointLock extends DisjointLock{
  /** The number of threads of the two types using the resource.  Invariant: 
    * £aIn = 0£ or £bIn = 0£. */
  private var aIn, bIn = 0

  /** The number of threads of the two types waiting. */
  private var aWaiting, bWaiting = 0

  private val lock = new Lock

  /** Conditions for signalling to the two types of threads.  A signal on
    * £aSignal£ indicates £bIn = 0£, and similarly with £bSignal£. */
  private val aSignal, bSignal = lock.newCondition

  def aEnter() = lock.mutex{
    if(bIn > 0 || bWaiting > 0){ // Thread has to wait.
      aWaiting += 1
      aSignal.await(bIn == 0) // Wait for B£-£threads to leave (1).
      aWaiting -= 1
    }
    aIn += 1
  }

  def aExit() = lock.mutex{ 
    aIn -= 1
    if(aIn == 0) bSignal.signalAll() // Signal to B£-£threads at (2).
  }

  def bEnter() = lock.mutex{
    if(aIn > 0 || aWaiting > 0){ // Thread has to wait.
      bWaiting += 1
      bSignal.await(aIn == 0) // Wait for A£-£threads to leave (2).
      bWaiting -= 1
    }
    bIn += 1
  }

  def bExit() = lock.mutex{
    bIn -= 1
    if(bIn == 0) aSignal.signalAll() // Signal to A£-£threads at (1).
  }
}
\end{scala}

Note that it is necessary for a thread to recheck the relevant condition when
it receives a signal.  Consider the following scenario.
%
\begin{enumerate}
\item[1.] An A-thread is using the resource; two B-threads are waiting at (2).
\item[2.] The A-thread exists and signals to both B-threads.
\item[3.] One B-thread is scheduled and enters the critical section, but the
  other is not yet scheduled.
\item[4.] Another A-thread calls |aEnter| and waits at (1).
\item[5.] The first B-thread exits and signals to the A-thread.
\item[6.] The A-thread enters.
\item[7.] The second B-thread is scheduled.
\end{enumerate}
%
At this point, the second B-thread needs to recheck the condition: if it
simply enters the critical region, it breaks the intended invariant. 

We can test this lock using linearizability testing, where the abstract state
is a pair $(a,b)$ indicating $a$ A-threads and $b$ B-threads are using the
resource, with the invariant $a = 0$ or $b = 0$.  The corresponding sequential
operations are straightforward; those for entering capture the obvious
precondition.
%
\begin{scala}
  type S = (Int, Int)
  def seqAEnter(s: S): (Unit, S) = {
    val (a,b) = s; require(b == 0); ((), (a+1,b))
  }
  def seqAExit(s: S): (Unit, S) = {
    val (a,b) = s; assert(a > 0 && b == 0); ((), (a-1,b))
  }
  def seqBEnter(s: S): (Unit, S) = {
    val (a,b) = s; require(a == 0); ((), (a,b+1))
  }
  def seqBExit(s: S): (Unit, S) = {
    val (a,b) = s; assert(b > 0 && a == 0); ((), (a,b-1))
  }
\end{scala}
%
We can then run threads that alternate between entering and leaving, as a
particular type of thread.  The rest of the linearizability testing is
standard. 
\end{answerS}
